\section{Многочлены степени не выше трёх}

Для многочленов малых степеней транзитивные подгруппы $S_n$ легко перечислить, и группа Галуа определяется
тем, приводим ли многочлен и является ли его дискриминант квадратом.

\begin{example}[4.6]
Пусть $f=X^2+bX+c\in F[X]$, $\operatorname{char}F\ne2$, и $f$ не является квадратом линейного многочлена, т.\,е.\
$D(f)=b^2-4c\ne0$ ($f$ сепарабелен). Тогда
\[
f\ \text{неприводим}\iff D(f)\ \text{не квадрат в }F\iff G_f=S_2 .
\]
Действительно, $f=\bigl(X+\frac b2\bigr)^2-\frac{b^2-4c}4$, корни $f$ равны $\frac{-b\pm\sqrt{D}}2$, и многочлен степени $2$
неприводим ровно тогда, когда у него нет корней в $F$, т.\,е.\ когда $\sqrt D\notin F$. Вторая равносильность~--- это
\rf{4.2} ($A_2=1$), или \rf{4.5}: $S_2$~--- единственная транзитивная подгруппа $S_2$.

В характеристике $2$ оба условия <<портятся>>.
\begin{enumerate}[label=(\alph*)]
\item Многочлен может быть неприводимым, но несепарабельным: $X^2-a$ при $a\in F\smallsetminus F^2$ (например, $F=\F_2(t)$,
$a=t$); в поле разложения $X^2-a=(X-\sqrt a)^2$.
\item Многочлен может быть неприводимым и сепарабельным, но с дискриминантом-квадратом. Пусть $f=X^2-X-a=X^2+X+a$, и
$\alpha$~--- его корень; второй корень~--- $\alpha+1$ (сумма корней равна $1$), так что $f=(X-\alpha)(X-\alpha-1)$ и
$D(f)=\bigl(\alpha-(\alpha+1)\bigr)^2=1$~--- квадрат при любом $a$. Но такой $f$ может быть неприводимым:
например, $X^2+X+1$ над $\F_2$ (\rf{4.4}); тогда $G_f=S_2\not\subset A_2$.
\end{enumerate}
Для $f=X^2+X+a$ дискриминант Берлекэмпа равен $a$ (вычисление~--- в конце \rf{4.4}), и критерий \rf{4.4}
говорит: $G_f\subset A_2=1$ ($f$ раскладывается над $F$) $\iff$ $a=\delta^2+\delta$ для некоторого $\delta\in F$~--- что очевидно и
напрямую ($\delta$~--- корень $f$). К этому случаю сводится любой сепарабельный $X^2+bX+c$ в характеристике $2$: сепарабельность
означает $b\ne0$, и замена $X=bY$ даёт $b^2\bigl(Y^2+Y+\frac c{b^2}\bigr)$ (многочлен Артина--Шрайера, \rf{2.22b}).
\end{example}

\begin{example}[4.7]
Пусть $f\in F[X]$~--- многочлен степени $3$ и $\operatorname{char}F\ne2,3$. Если $f$ приводим, то он раскладывается на
линейный множитель и многочлен степени $\le2$, и мы по существу возвращаемся к \rf{4.6} (см.\ \rf{4.7a}). Пусть $f$ неприводим.
Тогда $f$ сепарабелен: неприводимый многочлен с кратным корнем в характеристике $p>0$~--- многочлен от $X^p$ (\rf{2.20}), и его
степень делится на $p$, а $3$ делится только на $p=3$. По \rf{4.5} $G_f$~--- транзитивная подгруппа $S_3$, и её порядок
делится на $3$ (замечание после \rf{4.5}). Таких подгрупп две: $A_3=\langle(1\,2\,3)\rangle$ и $S_3$. По \rf{4.2}
\[
G_f=A_3\iff D(f)\ \text{квадрат в }F,\qquad G_f=S_3\iff D(f)\ \text{не квадрат в }F .
\]
Например, $X^3-3X+1\in\Q[X]$ неприводим (это $-h(-X)$ для многочлена $h=X^3-3X-1$ из \rf{1.12}; или прямо по \rf{1.11}:
$\pm1$ не корни), его дискриминант $-4(-3)^3-27\cdot1^2=81=9^2$, и его группа Галуа~--- $A_3$
(ср.\ \rf{3.22a}, где это поле найдено явно: $\Q[2\cos\frac{2\pi}9]$). Многочлен $X^3+3X+1\in\Q[X]$ тоже неприводим (\rf{1.11}), но его
дискриминант $-4\cdot3^3-27=-135$ не квадрат, и группа Галуа~--- $S_3$ (\rf{3.23c}).

Условие $\operatorname{char}F\ne2$ здесь существенно: в характеристике $2$ дискриминант всегда квадрат (\rf{4.4}). Например, пусть
$F=\F_2(t)$ и $f=X^3-t$. У $f$ нет корня в $F$ (из $(u/v)^3=t$ с $u,v\in\F_2[t]$ следовало бы $3(\deg u-\deg v)=1$), так что
кубический многочлен $f$ неприводим; он сепарабелен ($f'=3X^2=X^2$ и $f$ взаимно просты). Если $\tau$~--- корень $f$, то корни~---
$\tau,\omega\tau,\omega^2\tau$, где $\omega$~--- корень $X^2+X+1$, и поле разложения содержит $\omega=\omega\tau/\tau$. Многочлен $X^2+X+1$
не имеет корней в $F$ (корень $u/v$ с взаимно простыми $u,v\in\F_2[t]$ должен иметь $u,v\mid1$, т.\,е.\ лежать в $\F_2$, а в
$\F_2$ корней нет), поэтому степень поля разложения делится на $2$ и на $3$, и $G_f=S_3$. Тем не менее $D(f)=-27t^2=t^2$~---
квадрат.
\end{example}

\begin{remark}[4.7a: сводка для кубических многочленов]
Пусть $\operatorname{char}F\ne2$ и $f\in F[X]$~--- сепарабельный многочлен степени $3$. Группа $G_f\subset S_3$ определяется
числом корней $f$ в $F$ и тем, является ли $D(f)$ квадратом:
\begin{center}
\begin{tabular}{@{}lll@{}}
\toprule
корней в $F$ & $D(f)$ & $G_f$\\
\midrule
$3$ & квадрат & $1$\\
$1$ & не квадрат & $S_2=\langle(2\,3)\rangle$ (транспозиция)\\
$0$ & квадрат & $A_3$\\
$0$ & не квадрат & $S_3$\\
\bottomrule
\end{tabular}
\end{center}
Ровно два корня в $F$ у кубического многочлена быть не может (третий равен $-a-x_1-x_2$, если $f=X^3+aX^2+\dots$). Если
$f=(X-x_1)h$ с неприводимым $h$, то $G_f=G_h=S_2$ (\rf{3.27e}) и $D(f)=D(h)\,h(x_1)^2$ не квадрат, так как $D(h)$ не квадрат
(\rf{4.6}). Над $F\subset\R$ знак дискриминанта показывает число вещественных корней (\rf{4.3}): при $D(f)<0$ вещественный корень
один, и $G_f\not\subset A_3$; при $D(f)>0$ все три корня вещественны.

\emph{Два семейства над $\Q$.} Семейство циклических кубик~--- \emph{многочлены Шенкса}
\[
S_a=X^3-aX^2-(a+3)X-1,\qquad D(S_a)=(a^2+3a+9)^2\quad\text{(по 4.2a)}.
\]
При $a\in\Z$ многочлен $S_a$ неприводим: рациональные корни унитарного многочлена с целыми коэффициентами и свободным
членом $-1$ могут быть только $\pm1$ (\rf{1.11}), а $S_a(1)=-2a-3\ne0$, $S_a(-1)=1$. Так как $D(S_a)$~--- квадрат, $G_{S_a}=A_3$ (\rf{4.7}).
Образующую этой группы можно указать явно: если $\rho$~--- корень $S_a$, то прямая проверка даёт $(1+\rho)^3S_a\bigl(-\frac1{1+\rho}\bigr)=-S_a(\rho)=0$, т.\,е.\
$\rho'=-\frac1{1+\rho}\in\Q[\rho]$~--- тоже корень ($\rho\ne-1$, так как $S_a(-1)\ne0$). Он отличен от $\rho$ (иначе $\rho^2+\rho+1=0$, а
$\rho$ имеет степень $3$), поэтому автоморфизм $\rho\mapsto-\frac1{1+\rho}$ поля $\Q[\rho]$ нетривиален и порождает $G_{S_a}\cong C_3$. При
$a=-1$ получаем $X^3+X^2-2X-1$~--- минимальный многочлен числа $2\cos\frac{2\pi}7$ (\rf{3.22}; $D=49$), при $a=0$~--- многочлен
$X^3-3X-1$ из \rf{1.12} ($D=81$).

<<Типичная>> же кубика имеет группу $S_3$. При $pq\ne0$ замена $X=\frac qpY$ даёт
\[
X^3+pX+q=\frac{q^3}{p^3}\bigl(Y^3+AY+A\bigr),\qquad A=\frac{p^3}{q^2},
\]
так что поле разложения $X^3+pX+q$ зависит только от $A$. Для переменной $A$ (над полем $\Q(A)$) многочлен $Y^3+AY+A$
неприводим по критерию Эйзенштейна для простого элемента $A$ кольца $\Q[A]$ (ср.\ \rf{1.16d}), а $D(Y^3+AY+A)=-A^2(4A+27)$ не
квадрат в $\Q(A)$ (многочлен $-(4A+27)$ имеет нечётную степень); по \rf{4.7} его группа Галуа над $\Q(A)$ равна $S_3$.
\end{remark}

\subsection*{Формула Кардано: дополнения}

Вывод формулы Кардано для $X^3+pX+q$ приведён в \rf{3.28a}: ищем корень в виде $x=u+v$ и \emph{дополнительно} требуем
$3uv=-p$; тогда $u^3+v^3=-q$, и $u^3,v^3$~--- корни квадратного уравнения $Z^2+qZ-\frac{p^3}{27}=0$. Наложив лишнее условие,
мы могли, казалось бы, потерять часть корней. Покажем, что этого не происходит, и уточним выбор кубических корней.

\begin{theorem}[4.7b: формула Кардано]
Пусть $\operatorname{char}F\ne2,3$, $f=X^3+pX+q\in F[X]$, $\Om\supset F$~--- алгебраически замкнутое поле и $\omega\in\Om$~---
первообразный кубический корень из $1$. Пусть $z_1,z_2\in\Om$~--- корни многочлена $Z^2+qZ-\frac{p^3}{27}$, занумерованные так,
что $z_1\ne0$ (если $p=q=0$, то $f=X^3$ и доказывать нечего). Выберем $u\in\Om$ с $u^3=z_1$ и положим $v=-\frac p{3u}$. Тогда
$v^3=z_2$ и
\[
f=\bigl(X-(u+v)\bigr)\bigl(X-(\omega u+\omega^2v)\bigr)\bigl(X-(\omega^2u+\omega v)\bigr),
\]
т.\,е.\ три числа $u+v$, $\omega u+\omega^2v$, $\omega^2u+\omega v$~--- это \emph{все} корни $f$ с учётом кратности.
\end{theorem}

\begin{proof}
Так как $z_1z_2=-\frac{p^3}{27}$, имеем $v^3=-\frac{p^3}{27z_1}=z_2$ (при $p=0$ это верно, поскольку тогда $z_2=0=v$). Далее
$3uv=-p$ и $u^3+v^3=z_1+z_2=-q$, и тождество из \rf{3.28a} (прямое перемножение с учётом $1+\omega+\omega^2=0$)
\[
\bigl(X-(u+v)\bigr)\bigl(X-(\omega u+\omega^2v)\bigr)\bigl(X-(\omega^2u+\omega v)\bigr)=X^3-3uv\,X-(u^3+v^3)
\]
даёт требуемое.
\end{proof}

\begin{remark}[4.7c]
(a) \emph{Почему дополнительное условие $3uv=-p$ не теряет корней.} Теорема \rf{4.7b} уже даёт ответ: при этом условии три
числа $\omega^ku+\omega^{-k}v$ исчерпывают корни. Другое объяснение: \emph{любой} корень $x$ многочлена $f$ можно записать в
виде $x=u+v$ с $3uv=-p$~--- достаточно взять в качестве $u,v$ корни многочлена $T^2-xT-\frac p3$ (они есть в $\Om$). Тогда
$u^3+v^3=(u+v)^3-3uv(u+v)=x^3+px=-q$. Значит, условие $3uv=-p$ не сужает множество корней, а лишь фиксирует, \emph{как}
корень разбит на два слагаемых.

(b) \emph{Согласование кубических корней.} В записи
$x=\sqrt[3]{-\frac q2+\sqrt{\frac{q^2}4+\frac{p^3}{27}}}+\sqrt[3]{-\frac q2-\sqrt{\frac{q^2}4+\frac{p^3}{27}}}$ каждый кубический корень имеет три
значения, т.\,е.\ формально получается девять сумм $\omega^ju+\omega^kv$. Если $pq\ne0$, корнями $f$ являются только три из
них~--- те, где $j+k\equiv0\pmod 3$, т.\,е.\ слагаемые $u'=\omega^ju$, $v'=\omega^kv$ удовлетворяют условию $u'v'=-\frac p3$. Действительно,
в общем случае $3u'v'=-\omega^{j+k}p$ и $u'^3+v'^3=-q$, и по тому же тождеству $u'+v'$~--- корень многочлена $X^3+\omega^{j+k}pX+q$. При
$j+k\not\equiv0$ общий корень $x$ этого многочлена и $f$ удовлетворял бы равенству $(\omega^{j+k}-1)px=0$, т.\,е.\ $x=0$ и $q=0$.
Замена $u$ на $\omega u$ (с $v$, заменённым на $\omega^2v$) лишь переставляет три корня. Вырожденные случаи: при $p=0$ имеем $v=0$, и все девять сумм~---
корни $f=X^3+q$; при $q=0$ (и $p\ne0$) одна из трёх <<законных>> сумм равна $0$ (так как $0$~--- корень $f$), скажем $\omega^ju+\omega^{-j}v=0$, и тогда
$\omega^{j+1}u+\omega^{1-j}v$ и $\omega^{j+2}u+\omega^{2-j}v$, т.\,е.\ две <<лишние>> суммы, тоже равны $0$~--- корню $f$. Например, для $X^3-X$
при $u=\frac1{\sqrt3}e^{i\pi/6}$, $v=\bar u$ <<лишняя>> сумма $u+\omega^2v$ равна $0$.

(c) \emph{Подстановка Виета.} Ту же формулу даёт одна неизвестная: при $x=w-\frac p{3w}$
\[
w^3f\Bigl(w-\frac p{3w}\Bigr)=w^6+qw^3-\frac{p^3}{27},
\]
т.\,е.\ $w^3$~--- корень $Z^2+qZ-\frac{p^3}{27}$, и $x=w+v$ с $v=-\frac p{3w}$, как в \rf{4.7b}.
\end{remark}

\subsection*{Почему формула Кардано работает: резольвенты Лагранжа}

\begin{proposition}[4.7d]
Пусть $\operatorname{char}F\ne2,3$, $f=X^3+pX+q\in F[X]$ сепарабелен, $x_1,x_2,x_3$~--- его корни в $\Om$, $\omega\in\Om$~---
первообразный кубический корень из $1$, $\sqrt{-3}:=\omega-\omega^2=2\omega+1$ и $\Delta=(x_1-x_2)(x_1-x_3)(x_2-x_3)$. Положим
\[
r_1=x_1+\omega x_2+\omega^2x_3,\qquad r_2=x_1+\omega^2x_2+\omega x_3
\]
(\emph{резольвенты Лагранжа}). Тогда
\[
r_1r_2=-3p,\qquad r_1^3+r_2^3=-27q,\qquad r_1^3-r_2^3=3\sqrt{-3}\,\Delta,
\]
и корни восстанавливаются по формулам
\[
x_1=\tfrac13(r_1+r_2),\qquad x_2=\tfrac13(\omega^2r_1+\omega r_2),\qquad x_3=\tfrac13(\omega r_1+\omega^2r_2).
\]
В частности, $u=\frac{r_1}3$ и $v=\frac{r_2}3$~--- числа из формулы Кардано, и
\[
u^3=-\frac q2+\frac{\sqrt{-3}}{18}\,\Delta,\qquad v^3=-\frac q2-\frac{\sqrt{-3}}{18}\,\Delta .
\]
\end{proposition}

\begin{proof}
Все равенства~--- тождества между многочленами от $x_1,x_2$ (при $x_3=-x_1-x_2$, так как $x_1+x_2+x_3=0$) с
$p=x_1x_2+x_1x_3+x_2x_3$, $q=-x_1x_2x_3$ и $\omega^2+\omega+1=0$; они проверяются раскрытием скобок. Например,
$r_1r_2=x_1^2+x_2^2+x_3^2+(\omega+\omega^2)(x_1x_2+x_1x_3+x_2x_3)=(x_1+x_2+x_3)^2-3p=-3p$, и
$\omega^2r_1+\omega r_2=(\omega^2+\omega)x_1+2x_2+(\omega+\omega^2)x_3=-x_1+2x_2-x_3=3x_2$. Формулы для $u^3,v^3$ получаются из
$u^3=\frac{(r_1^3+r_2^3)+(r_1^3-r_2^3)}{54}$; заметим, что $\bigl(\frac{\sqrt{-3}}{18}\Delta\bigr)^2=-\frac{D(f)}{108}=\frac{q^2}4+\frac{p^3}{27}$,
в согласии с \rf{3.28a}.
\end{proof}

\begin{remark}[4.7e: башня Кардано и группа $S_3$]
Пусть $E$~--- поле разложения $f$ над $F$ и $G=\Gal(E[\omega]/F[\omega])\subset G_f\subset S_3$. Цикл $(1\,2\,3)$ ($x_1\mapsto x_2\mapsto
x_3\mapsto x_1$) переводит $r_1$ в $x_2+\omega x_3+\omega^2x_1=\omega^2r_1$, а транспозиция $(2\,3)$ меняет $r_1$ и $r_2$ местами. Значит,
$r_1^3$ неподвижен относительно $G\cap A_3$ и лежит в $F[\omega,\Delta]$ (\rf{4.2}, применённое над $F[\omega]$; это видно и из
формулы для $u^3$ в \rf{4.7d}). Получаем башню
\[
F\ \subset\ F[\omega]\ \subset\ F[\omega,\Delta]\ \subset\ F[\omega,\Delta,r_1]=E[\omega],
\]
шаги которой~--- присоединение корня из $1$, квадратного корня из $D(f)$ и кубического корня из $r_1^3$ (последнее
равенство~--- из формул \rf{4.7d} для $x_i$, если $r_1\ne0$; если $r_1=0$, поменяем ролями $r_1$ и $r_2$). Над $F[\omega]$ это
цепочка подгрупп $S_3\supset A_3\supset1$ с циклическими факторами~--- пример общей теоремы \rf{3.28} о разрешимости в
радикалах. Поле $E[\omega]$ может быть больше $E$: для $X^3+3X+1$ радикалы Кардано лежат вне поля разложения (\rf{3.28a}).
\end{remark}

\begin{example}[4.7f]
(a) (Кардано, 1545) $X^3+6X-20$: $p=6$, $q=-20$, $\frac{q^2}4+\frac{p^3}{27}=108$, $z_{1,2}=10\pm6\sqrt3$. Так как
$(1\pm\sqrt3)^3=1\pm3\sqrt3+9\pm3\sqrt3=10\pm6\sqrt3$, можно взять $u=1+\sqrt3$; тогда $v=-\frac{2}{1+\sqrt3}=1-\sqrt3$, и корни равны
$u+v=2$, $\omega u+\omega^2v=-1+\sqrt3(\omega-\omega^2)=-1+3i$ и $-1-3i$; действительно, $X^3+6X-20=(X-2)(X^2+2X+10)$. Формула
<<прячет>> рациональный корень за выражением $\sqrt[3]{10+6\sqrt3}+\sqrt[3]{10-6\sqrt3}$.

(b) (Бомбелли, 1572) $X^3-15X-4$: $p=-15$, $q=-4$, $\frac{q^2}4+\frac{p^3}{27}=4-125=-121$, $z_{1,2}=2\pm11i$. Так как
$(2\pm i)^3=2\pm11i$, берём $u=2+i$, $v=\frac5{2+i}=2-i$, и корни равны $u+v=4$ и
$\omega u+\omega^2v=2(\omega+\omega^2)+i(\omega-\omega^2)=-2-\sqrt3$, $\omega^2u+\omega v=-2+\sqrt3$~--- все вещественны. Здесь
$D=13068>0$, и потому $\frac{q^2}4+\frac{p^3}{27}=-\frac D{108}<0$: при трёх вещественных корнях формула Кардано \emph{всегда}
проходит через невещественные числа. Многочлен $X^3-15X-4=(X-4)(X^2+4X+1)$ приводим, и кубические корни удаётся
<<распутать>>; для неприводимого многочлена с тремя вещественными корнями никакая запись вещественными радикалами
невозможна (\rf{3.22f}).
\end{example}

\subsection*{Тригонометрическое решение}

В неприводимом случае ($D>0$) кубический корень из комплексного числа $z_1$ сводится к делению его аргумента на три.

\begin{proposition}[4.7g]
Пусть $p,q\in\R$ и $D=-4p^3-27q^2>0$. Тогда $p<0$, число $\beta=\frac{3q}{2p}\sqrt{-\frac3p}$ лежит в интервале $(-1,1)$, и при
$\beta=\cos\varphi$, $0<\varphi<\pi$, корни многочлена $X^3+pX+q$ равны
\[
x_k=2\sqrt{-\tfrac p3}\,\cos\frac{\varphi+2\pi k}3,\qquad k=0,1,2 .
\]
\end{proposition}

\begin{proof}
Из $-4p^3>27q^2\ge0$ следует $p<0$. Положим $\kappa=2\sqrt{-p/3}$, так что $p=-\frac34\kappa^2$, и $x=\kappa y$. Тогда
\[
x^3+px+q=\frac{\kappa^3}4\bigl(4y^3-3y\bigr)+q,
\]
и уравнение $f(x)=0$ равносильно $4y^3-3y=-\frac{4q}{\kappa^3}=\beta$ (прямая проверка). Далее $\beta^2=-\frac{27q^2}{4p^3}$, откуда
$1-\beta^2=\frac{D}{-4p^3}>0$. При $y=\cos\theta$ имеем $4y^3-3y=\cos3\theta$, поэтому числа $y_k=\cos\frac{\varphi+2\pi k}3$ ($k=0,1,2$)~---
корни $4y^3-3y-\beta$. Они различны: углы $\frac\varphi3$, $\frac{\varphi+2\pi}3$, $\frac{\varphi+4\pi}3$ лежат в интервалах $\bigl(0,\frac\pi3\bigr)$,
$\bigl(\frac{2\pi}3,\pi\bigr)$, $\bigl(\frac{4\pi}3,\frac{5\pi}3\bigr)$, и косинусы на них принимают значения из $\bigl(\frac12,1\bigr)$, $\bigl(-1,-\frac12\bigr)$,
$\bigl(-\frac12,\frac12\bigr)$ соответственно. У многочлена степени $3$ других корней нет. (Ср.\ \rf{1.42a}.)
\end{proof}

\begin{remark}[4.7h]
(a) \emph{Связь с формулой Кардано.} При $D>0$ числа $z_{1,2}=-\frac q2\pm i\sqrt{D/108}$ комплексно сопряжены, и
$|z_1|^2=z_1z_2=-\frac{p^3}{27}$, т.\,е.\ $|z_1|=\bigl(-\frac p3\bigr)^{3/2}$ и $z_1=|z_1|e^{i\varphi}$ с тем же $\varphi$
($\cos\varphi=-\frac q2|z_1|^{-1}=\beta$). Кубический корень $u=\sqrt{-p/3}\,e^{i\varphi/3}$ получается извлечением \emph{квадратного} корня
из $-\frac p3$ и делением угла на три, $v=\bar u$, и $u+\bar u=2\sqrt{-p/3}\cos\frac\varphi3=x_0$. Таким образом, корни кубического
многочлена с $D>0$ строятся циркулем, линейкой и \emph{трисектором} (инструментом для деления угла на три части).
При $D<0$ трисектор не помогает: рассуждение \rf{1.42a}, доказывающее, что $\sqrt[3]2$ им не построить, дословно
переносится на вещественный корень любого неприводимого над $\Q$ кубического многочлена с $D<0$ (этот корень~--- единственный
вещественный).

(b) Для $X^3-3X+1$: $\beta=-\frac12$, $\varphi=\frac{2\pi}3$, и корни~--- $2\cos\frac{2\pi}9$, $2\cos\frac{8\pi}9$, $2\cos\frac{14\pi}9=2\cos\frac{4\pi}9$, как
в \rf{3.22a}. Для $X^3-15X-4$: $\beta=\frac2{5\sqrt5}$, и $x_0=2\sqrt5\cos\frac\varphi3=4$, так как $\cos\frac\varphi3=\frac2{\sqrt5}$ (проверка:
$4\bigl(\frac2{\sqrt5}\bigr)^3-3\cdot\frac2{\sqrt5}=\frac2{5\sqrt5}$).

(c) При $D<0$ вещественный корень один, и вещественная формула Кардано выражает его вещественными радикалами
(\rf{3.23c}, \rf{3.28a}).
\end{remark}

\begin{aside}[4.7i: преобразование Чирнгауза]
Другой путь к решению: квадратичная замена $y=x^2+\alpha x+\beta$ переводит корни $x_i$ многочлена $X^3+pX+q$ в числа
$y_i=x_i^2+\alpha x_i+\beta$, и параметры можно подобрать так, чтобы $y_1+y_2+y_3=0$ и $y_1^2+y_2^2+y_3^2=0$; тогда
$y_1y_2+y_1y_3+y_2y_3=0$, и $y_i$~--- корни уравнения $Y^3=N$, $N=y_1y_2y_3\in F[\alpha]$. Используя $\sum x_i=0$, $\sum x_i^2=-2p$,
получаем $\sum y_i=-2p+3\beta$, т.\,е.\ $\beta=\frac{2p}3$, а затем (прямое вычисление)
\[
\textstyle\sum y_i^2=-\frac23\bigl(3p\alpha^2+9q\alpha-p^2\bigr).
\]
Дискриминант квадратного уравнения $3p\alpha^2+9q\alpha-p^2=0$ равен $81q^2+12p^3=-3D$. Снова сначала извлекается квадратный
корень (из $-3D$), затем кубический (из $N$)~--- та же структура, что в \rf{4.7e}.
\end{aside}
